% !TEX root = ../main_zh.tex
\chapter{人口、土地与马尔萨斯经济}
\label{ch:malthus}

第~\ref{ch:solow}~章的索罗模型，以及第~\ref{ch:neoclassical}~章为它补上微观基础的新古典版本，都暗含两个未曾言明的承诺：它们只用资本与劳动来构造产出；它们把人口规模当成从外部交给经济的东西——一个由建模者随手写下、以速率 $n$ 增长的数字。对于这两个模型本来要解释的世界——一个紧约束在于可累积的资本、人们越富裕反而生得\emph{越少}的现代经济——这两个承诺都站得住脚。可一旦面对它之前的那个世界，二者就都失灵了。在前工业的农业社会里，稀缺要素并不是资本，而是\emph{土地}：土地是固定的，无法被累积。人口也不是外生的：一个社会里究竟有多少人，本身就是一项经济决策的结果——是各个家庭关于养育多少孩子的决策。

本章研究的，正是把这两件事重新放回中心的模型。它建立在一个一旦看清便难以再忽视的典型事实之上：纵观漫长的农业时代，技术进步并没有让普通人过得更好。更好的种子、更好的犁、更好的灌溉，最终没有体现为更高的生活水平，而是体现为\emph{更多的人}以和从前一样的水平活着。其中的机制正是托马斯·马尔萨斯（Thomas Malthus）在 1798 年描述的那一套：当某项技术改良提高了土地的产量，家庭会短暂地变富，于是多养孩子；人口随之增长；固定的那块土地被更多张嘴瓜分；人均产量又被压回到原来的位置。土地是紧约束，而人口正是这条约束所要按住的那个变量。整个模型的用意，就是把这个陷阱说精确——然后再追问：人类最终究竟靠什么逃出了它。

我们分三步走。先把土地重新放回生产函数，说明为什么一个固定要素会把日益增加的投入变成不断递减的回报——这就是农业的\emph{过密化（involution）}现象。接着写出并求解家庭的最优化问题，导出人口的运动方程，定位出马尔萨斯陷阱真正咬住经济的那个稳态：生活水平被钉死在维生水平，任何技术改进都被全数花在了多出来的人口上。最后，我们把养育孩子的成本修改为还包含一项\emph{时间}成本，它会把收入与生育之间的关系压平、从而拆掉这个陷阱；再叠加一个数量—质量的取舍，把这条关系彻底翻转过来——这便是人口转变。

\section{土地作为紧约束要素}

我们目前见过的增长模型，其生产函数里只有资本 $K$ 和劳动 $L$ 两个要素。要刻画一个农业经济，我们再加入第三个要素——土地 $T$，并写下一个柯布—道格拉斯技术
\begin{equation}
F(K, L, T) = A\, K^{\alpha} L^{\beta} T^{\,1-\alpha-\beta},
\label{eq:malthus-prod}
\end{equation}
其中各指数均为正、且相加等于一，因此该技术对三种要素整体而言是规模报酬不变的。$A$ 是全要素生产率，代表技术的水平。

在\emph{现代}经济里，土地几乎无足轻重。相对于生产所需，土地实在太多，土地约束从来咬不住，产出由资本和劳动主宰，于是土地份额 $1-\alpha-\beta$ 可以略去而几乎没有损失——这恰恰是第~\ref{ch:solow}~章和第~\ref{ch:neoclassical}~章把它丢掉的原因。可在农业社会里情形正好相反。那里可耕地的供给实际上是固定的，一旦不断增长的人口把它全都开垦殆尽，土地就变成了\emph{硬约束}。此时往一块固定的地上不断追加劳动力，就会狠狠撞上边际报酬递减规律：每多一双手为总产出带来的增量越来越小，到最后劳动的边际产出远远跌到了平均产出之下。

\defn{农业过密化}{
\emph{过密化}（亦即“高水平均衡陷阱”）指的是这样一种局面：一个固定要素——此处即土地——逼着人们投入越来越多的劳动，却只换来越来越小的边际产出增量。人们把心力与巧思一股脑倾注进去，只为从同一块地里多榨出一点点；总产出甚至可能还在上升，但追加劳动的边际产出已坍缩，\emph{人均}产出陷入停滞。那种极尽精细、劳力饱和的耕作方式，正是一个社会被死死压在土地约束上的可见症状。
}

\rmkb[一个熟悉模式的名字]{“过密化”（involution，中文译作\emph{内卷}）这个词近来已经跑出了农业史的研讨课，被用来形容上班族为同一批固定的奖赏越拼越狠。其经济骨架是完全一样的：一个固定要素、不断加码的努力，以及一项已被压到几乎为零的努力边际回报。马尔萨斯模型正是这一模式最早、最农业的那个版本。}

至此我们可以陈述这个模型存在的理由——它要解释的那个核心经验规律。

\state{马尔萨斯典型事实}{
在农业社会里，人们很难看到任何持续的人均生活水平的改进。技术进步带来的，反而是一个在生活水平不变的前提下\emph{规模更大的人口}。一个具有这种性质的经济——更好的技术所带来的收益被更多的人、而非更富的人吸收掉——就被称为\emph{马尔萨斯（Malthus）经济}。
}

本章余下的任务，就是把这个事实从最优化行为中推导出来，让我们不只是知道它会发生，而是确切地理解它为什么发生。

\section{家庭的最优化问题}

考虑农业经济中的一个代表性家庭。它在乎自己的消费 $C_t$，也在乎自己养育的孩子数量 $n_t$，并按速率 $\beta \in (0,1)$ 对未来进行贴现。其终生效用为
\begin{equation}
\max_{\{C_t,\, n_t\}_{t=0}^{\infty}} \;\sum_{t=0}^{\infty} \beta^{t}\pr{\log C_t + \log n_t}.
\label{eq:malthus-obj}
\end{equation}
对数设定承载了两层值得停下来体会的经济含义。孩子直接进入效用，刻画的是农业现实中的一个事实：孩子既是天伦之乐，也是一笔资产——田里多一双手、老来多一份保障。而效用取对数，意味着消费与孩子是弹性为一的不完全替代品，于是家庭总会两样都要一些，并按固定比例在二者间分配自己的资源——这一点我们马上就会看到。

每一期家庭都面对一个预算约束。养孩子是有代价的：每养一个要花掉 $b$ 单位的消费品。记 $c_t$ 为每个家庭的消费、$y_t$ 为每个家庭的收入，则花在消费和孩子上的资源不能超过收入，
\begin{equation}
c_t + b\, n_t = y_t,
\label{eq:malthus-budget}
\end{equation}
其中 $b > 0$ 是养育一个孩子的物质成本。

把这个家庭问题和上一节的土地约束系在一起的，是收入 $y_t$ 的生成方式。设 $P_t$ 为人口规模——它将是我们的\emph{状态变量}，即那个缓慢移动、概括了经济全部历史的存量。土地总量是固定的；把它标准化为一。那么人均土地就只是人口的倒数，
\begin{equation}
L_t = \frac{1}{P_t}.
\label{eq:malthus-land-per-person}
\end{equation}
人均产出取决于每个人能耕作多少土地。我们现在把生产函数~\eqref{eq:malthus-prod}~特化到以劳动与土地为起作用的要素的农业情形。为了让记号清楚，把劳动投入写成工人数 $N_t$（即先前称作 $L$ 的那个符号），把固定的土地存量写成 $T=1$，于是人均土地就是 $L_t = T/P_t = 1/P_t$，与~\eqref{eq:malthus-land-per-person}~一致——从这里起，字母 $L$ 表示人均土地，而不再表示劳动。把两要素产出 $F(N_t, T)$ 除以工人数，就得到一个关于人均土地递增且凹的人均产量，
\begin{equation}
y_t = \frac{F(N_t, T)}{N_t} = f(L_t) = z_t\, L_t^{\theta} = \frac{z_t}{P_t^{\theta}},
\qquad 0 < \theta < 1,
\label{eq:malthus-yield}
\end{equation}
其中 $z_t$ 是农业技术的水平（扮演上文 $A$ 的角色），$\theta$ 是产量对人均土地的弹性。最后一个等号用到了~\eqref{eq:malthus-land-per-person}~中的 $L_t = 1/P_t$。方程~\eqref{eq:malthus-yield}~是整个故事的发动机：\textbf{人口越多，人均土地越少，因而人均收入越低。}这正是那个固定要素伸进每一个家庭预算里的通道。

\asm{马尔萨斯经济}{
本模型依赖以下一组维持性假设：
\begin{itemize}
    \item 一个代表性家庭，其偏好~\eqref{eq:malthus-obj}~定义在消费 $C_t$ 与孩子数量 $n_t$ 之上，按 $\beta\in(0,1)$ 贴现。
    \item 每养一个孩子花费 $b$ 单位的消费品，得到预算约束 $c_t + b\,n_t = y_t$。
    \item 土地总量固定、标准化为一，故人均土地为 $L_t = 1/P_t$。
    \item 一个凹的人均产量 $y_t = z_t L_t^{\theta} = z_t / P_t^{\theta}$，其中 $0<\theta<1$，$z_t$ 为技术水平。
    \item 人口通过生育演化：下一期人口等于本期人口乘以每个家庭的孩子数。
\end{itemize}
}

\subsection{静态选择：如何分配收入}

家庭的完整问题是动态的，因为今天养育的孩子会扩大明天的人口，从而通过~\eqref{eq:malthus-yield}~压低明天的收入。但\emph{如何把给定的收入} $y_t$ 在消费和孩子之间分配，却是一个纯粹的期内静态问题：固定住 $y_t$，期内收益 $\log c_t + \log n_t$ 只取决于第 $t$ 期的选择。因此我们可以逐期求解这个分配。

\ex[期内分配]{
给定收入 $y_t$，选择 $c_t$ 与 $n_t$ 以最大化 $\log c_t + \log n_t$，约束为预算 $c_t + b\,n_t = y_t$。
}

\sol{
构造拉格朗日函数
\[
\mathcal{L} = \log c_t + \log n_t + \lambda\pr{y_t - c_t - b\, n_t}.
\]
一阶条件为
\[
\pd{\mathcal{L}}{c_t}: \quad \frac{1}{c_t} = \lambda,
\qquad
\pd{\mathcal{L}}{n_t}: \quad \frac{1}{n_t} = \lambda\, b.
\]
第一式除以第二式消去乘子，得到 $n_t b = c_t$：家庭花在养孩子上的钱，恰好和花在自己消费上的钱一样多。这正是对数—对数偏好的特征——每种用途都拿到一份固定的预算份额，此处即各占一半。把 $c_t = b\,n_t$ 代回预算 $c_t + b\,n_t = y_t$，得 $2 b\, n_t = y_t$，于是
\[
n_t^{*} = \frac{y_t}{2b}, \qquad c_t^{*} = \frac{y_t}{2}.
\]
}

这两条政策规则既干净又直观。收入的一半被消费掉；另一半花在养孩子上，按单位成本 $b$ 折算，可以买到 $y_t/(2b)$ 个孩子。关键在于，\textbf{生育随收入而上升}：越富的家庭养越多的孩子。仅凭这一条行为事实，再加上土地约束，就足以生成马尔萨斯陷阱。

\section{人口动态与稳态}

现在让静态规则去喂动态。每个家庭养 $n_t$ 个孩子，故下一期的人口为
\begin{equation}
P_{t+1} = n_t\, P_t.
\label{eq:malthus-lom}
\end{equation}
（把 $n_t$ 读作每个家庭的孩子数；在每个家庭一名成年人的设定下，它也就是整个人口的总繁殖系数。）把最优生育 $n_t^{*} = y_t/(2b)$、再把产量关系 $y_t = z_t / P_t^{\theta}$ 代入，就得到一个关于状态 $P_t$ 的单一差分方程，
\begin{equation}
P_{t+1} = \frac{y_t}{2b}\, P_t = \frac{z_t}{2b}\, P_t^{\,1-\theta}.
\label{eq:malthus-pdyn}
\end{equation}
指数 $1-\theta$ 严格介于零与一之间，故右端关于 $P_t$ 既凹又增：这个映射有唯一且稳定的正不动点。人口太少时会增长（土地充裕、收入高、家庭多生孩子）；人口太多时会收缩（土地稀缺、收入低、家庭少生孩子）。无论从哪一边出发，经济都被驱向同一个歇脚点。

\subsection{求解稳态}

在稳态上，人口不再变化，$P_{t+1} = P_t = P_{\mathrm{ss}}$，技术固定在 $z_t = z$。把这一条件加到~\eqref{eq:malthus-pdyn}~上，
\[
P_{\mathrm{ss}} = \frac{z}{2b}\, P_{\mathrm{ss}}^{\,1-\theta}
\quad\Longrightarrow\quad
P_{\mathrm{ss}}^{\,\theta} = \frac{z}{2b}
\quad\Longrightarrow\quad
P_{\mathrm{ss}} = \pr{\frac{z}{2b}}^{1/\theta}.
\]
其余一切都靠代入得到。由~\eqref{eq:malthus-yield}，稳态人均收入为
\[
y_{\mathrm{ss}} = \frac{z}{P_{\mathrm{ss}}^{\,\theta}}
= \frac{z}{\,z/(2b)\,} = 2b,
\]
再用最优规则便得到稳态的生育率与消费，
\[
n_{\mathrm{ss}} = \frac{y_{\mathrm{ss}}}{2b} = 1,
\qquad
c_{\mathrm{ss}} = \frac{y_{\mathrm{ss}}}{2} = b.
\]

\thm{马尔萨斯稳态}{
在土地存量固定、技术为 $z$、且面对上述家庭问题的情形下，经济收敛到唯一的稳态，其人口为
\[
P_{\mathrm{ss}} = \pr{\frac{z}{2b}}^{1/\theta},
\]
而在该稳态上，人均收入、生育率与人均消费分别为
\[
y_{\mathrm{ss}} = 2b, \qquad n_{\mathrm{ss}} = 1, \qquad c_{\mathrm{ss}} = b.
\]
}

这个稳态有两个特点值得着重强调。第一，$n_{\mathrm{ss}} = 1$：每个家庭恰好养育一个存活的孩子，于是人口刚好自我替代、停止增长。这正是对均衡最自然的读法——稳定的人口并非被假设进来，而是被经济学\emph{生产}出来的。当土地被瓜分到收入跌至 $y_{\mathrm{ss}} = 2b$ 时，各家恰好只养得起一个孩子，人口于是稳定下来。第二，也更引人注目：稳态的生活水平 $c_{\mathrm{ss}} = b$ 仅仅被钉在养育孩子的成本上。农业经济里的生活水平徘徊在一个由生物学与繁衍成本决定的位置，而不取决于这个社会技术有多高明。

\subsection{为何技术买来的是人口，而非繁荣}

这个模型的回报，正在于它对技术进步的论断。设技术水平 $z$ 永久上升——更好的种子、更重的犁、新的轮作制。看看那些稳态表达式，问问什么在动。

稳态人口上升：
\[
P_{\mathrm{ss}} = \pr{\frac{z}{2b}}^{1/\theta} \;\uparrow \qquad \text{当 } z\uparrow.
\]
但人均收入、生育率与人均消费却是
\[
y_{\mathrm{ss}} = 2b, \qquad n_{\mathrm{ss}} = 1, \qquad c_{\mathrm{ss}} = b,
\]
其中没有一个含 $z$。技术从每一个福利度量里都消失了。\textbf{更好的技术抬高了这块土地所能供养的人数，却抬不高其中任何一个人所享的生活水准。}这恰恰就是我们一开始要解释的那个典型事实，如今它是被推导出来的、而非被断言的：在马尔萨斯世界里，进步是以人口的形式支付出去的。

把这套机制顺着模型走一遍，可以干净地陈述成一条链。更高的 $z$ 在人口给定时抬高收入 $y_t$；更高的收入抬高生育 $n_t^{*} = y_t/(2b)$；更多的孩子扩大人口 $P_{t+1} = n_t P_t$；更大的人口意味着更少的人均土地 $L_t = 1/P_t$；而更少的人均土地又把收入 $y_t = z/P_t^{\theta}$ 压回去。经济就沿着这条链滑动，直到收入回到 $2b$、生育回到一——只不过此时人口更大了。更好的技术带来的那点暂时繁荣，到头来被完全转化成了多出来的嘴。

\rmkb[控制变量会跳，状态变量只爬]{这里值得留意不同变量的调整速度。生育与消费是\emph{控制}变量：家庭每一期都瞬间重置它们，所以当 $z$ 跳升时，收入和生育可以随之跳升。人口则是\emph{状态}变量，是一个只能通过出生与死亡缓慢变化的存量。于是一次技术改进会带来即刻的繁荣井喷与一波生育潮，随后是人口向新的、更高的稳态长久而缓慢的攀升——在这段攀升中，早先的繁荣被一点点侵蚀殆尽。好光景是真实的，却是过渡性的；多出来的人口则是永久的。}

\section{逃出陷阱：孩子的时间成本}

刚才搭起来的这个模型冷酷无情：每一步进步都被人口吞没，生活水平永远逃不出维生线。然而它们分明逃出去了——从工业化开始。这套机制里一定有什么东西断掉了。模型最重要的那个不足，恰恰指向了答案。

我们曾假设孩子的唯一成本就是物质成本 $b$。但养孩子还要花\emph{时间}——在一个劳动能挣到工资的经济里，那些时间本可以用来工作。因此一个孩子的真实成本，是花在他身上的物质\emph{加上}照料他时放弃的工资。设 $w_t$ 为工资，每个孩子占用一单位时间中的份额 $m$，则每个孩子的放弃收入为 $m\,w_t$。把预算按家庭的时间禀赋以工资计价重写，便成为
\begin{equation}
c_t + \pr{b + m\, w_t}\, n_t = w_t.
\label{eq:malthus-time-budget}
\end{equation}
右端是家庭把全部时间禀赋都用来工作时的价值；$(b + m w_t)$ 这一项则是一个孩子的\emph{完全}成本——物质加机会成本。这里决定性的新特征在于，这个完全成本\emph{随工资上升}：在更富的经济里，一个孩子所耗用的时间更值钱，于是孩子变得更贵。

像之前一样求解期内问题——对数—对数偏好再次把支出在消费与孩子之间对半分——得到
\[
c_t^{*} = \frac{w_t}{2},
\qquad
n_t^{*} = \frac{w_t}{2\pr{b + m\, w_t}}.
\]
现在追踪一下工资上升时会发生什么。消费无上限地上升，但生育不会。当 $w_t \to \infty$ 时，物质成本 $b$ 相对于时间成本变得微不足道，于是
\[
n_t^{*} = \frac{w_t}{2\pr{b + m w_t}} \longrightarrow \frac{1}{2m}.
\]
生育不再随收入失控暴涨；不断上升的孩子机会成本给它封了顶。那条曾经驱动马尔萨斯陷阱的失控反馈——越富意味着越多孩子、越多孩子意味着越穷——已在源头被拆除。随着工资继续上升，家庭愿意投在孩子单纯\emph{数量}上的资源份额不再增长，转而开始把收入替换到别的用途，包括对每个孩子投入更多。生育在收入上升时这样被压弯、封顶，便是人口转变。

\state{人口转变}{
一旦孩子既有物质成本又有时间成本，孩子的完全价格 $(b + m w)$ 便随工资上升。于是当经济变富时，生育不再无界地上升——它转而被压平，并渐近趋于 $1/(2m)$，而不是随收入暴涨。那条失控的马尔萨斯反馈被打断，人口会收敛，而不再永远跟着技术一路向上。富裕国家里观察到的生育率彻底\emph{下降}，则来自另一个更进一步的机制，即下文注释中展开的数量—质量取舍。这二者合在一起，就是\emph{人口转变（demographic transition）}。
}

\rmkb[为何时间成本只压弯生育，以及由谁来收尾]{朴素的对数—对数模型加上时间成本，给出的生育是 $n^{*}=w/[2(b+mw)]$，它其实关于工资是\emph{递增}的，只是上有 $1/(2m)$ 的界。所以时间成本本身只是移除了那条失控的马尔萨斯反馈、给生育封了顶；它并不能单独造出我们在富裕国家观察到的生育下降。真正带来彻底\emph{下降}的那一块，是数量—质量取舍。一旦家庭还能选择对\emph{每个}孩子投入多少——教育、健康、人力资本——更高的工资就会在边际上让质量更具吸引力，家庭于是从“多生、少投”替换向“少生、多投”。因此人口转变最好被理解为孩子的\emph{数量}与\emph{质量}之间的取舍：人力资本进入了生产函数，孩子的质量变得值得购买，被选择的生育率随之下降。}

同一套逻辑还解释了现代经济内部一个有充分记录的模式：越富的家庭、尤其是收入越高的女性，往往孩子越少。当脑力劳动取代体力劳动成为生产的关键投入、当女性工资上升时，孩子所需时间的机会成本随之攀升——而一旦数量—质量取舍让家庭得以替换向“投资于每个孩子”，被选择的孩子数量便下降。上升的时间成本压平了收入—生育关系；数量—质量取舍则把它倒转过来。那个孩子廉价、繁荣只会催生更多孩子的马尔萨斯世界，已被彻头彻尾地翻了过来。

\rmkb[关于人口政策的一点说明]{人们很容易把马尔萨斯模型读成压制人口增长的依据，二十世纪的人口控制项目有一部分理由正是这样讲的。但能为这种主张背书的，是那个\emph{农业}模型——其中产出被一个固定要素封顶，多一个人就是同一块地上多一张要喂的嘴。而一旦经济离开了农业——一旦它像索罗模型和新古典模型那样靠资本、想法和人力资本运转——这套逻辑就不再适用。人并不只是一块固定大饼的消费者；他们产生巨大的正外部性：发明、分工，并扩张着马尔萨斯模型当作常数按住的那个技术 $z$ 本身。把人口当成纯粹的负担，等于把马尔萨斯式的直觉硬塞进一个早已超越它的世界。这个模型是对农业过去的忠实描述，而不是对工业现在的处方。}

至此，土地与内生人口都已在手，我们也就走完了本课程长期、实物的那一面：索罗模型、它的微观基础版本，以及先于现代增长的那个马尔萨斯世界。接下来的几章将从实物经济转向名义经济——货币、银行与价格水平——从第~\ref{ch:money}~章开始。
